% ============================================================================
%  Template for a note. Copy it to content/notes/<name>.tex; it is published
%  at /notes/<name>/. It is also a normal LaTeX document: compile it with
%  pdflatex/latexmk as usual.
%
%  Supported on the web: sections (numbered), theorem-like environments,
%  equations with \label/\eqref, \ref, footnotes, lists, tables, figures
%  (use .png/.jpg/.svg images next to the .tex file), \cite with a .bib file,
%  and your own \newcommand's. TikZ pictures are not converted: export them
%  to .svg/.png and \includegraphics them.
% ============================================================================
\documentclass[11pt]{article}
\usepackage[utf8]{inputenc}
\usepackage{amsmath,amssymb,amsthm}
\usepackage{graphicx}
\usepackage{hyperref}

\newtheorem{theorem}{Theorem}
\newtheorem{lemma}[theorem]{Lemma}
\newtheorem{proposition}[theorem]{Proposition}
\newtheorem{corollary}[theorem]{Corollary}
\theoremstyle{definition}
\newtheorem{definition}[theorem]{Definition}
\newtheorem{example}[theorem]{Example}
\theoremstyle{remark}
\newtheorem{remark}[theorem]{Remark}

\newcommand{\E}{\mathbb{E}}
\newcommand{\R}{\mathbb{R}}

\title{Title of the note}
\author{Alberto Fernández-de-Marcos}
\date{September 2026}

\begin{document}
\maketitle

\begin{abstract}
One or two sentences describing the note.
\end{abstract}

\tableofcontents

\section{Introduction}\label{sec:intro}

\begin{definition}\label{def:sphere}
The unit hypersphere is $\mathbb{S}^q := \{\mathbf{x}\in\R^{q+1} : \|\mathbf{x}\| = 1\}$.
\end{definition}

\begin{theorem}[A named result]\label{thm:main}
If $X$ is uniform on $\mathbb{S}^q$, then $\E[X] = 0$.
\end{theorem}

\begin{proof}
By symmetry, $X$ and $-X$ have the same distribution, so
\begin{equation}\label{eq:sym}
  \E[X] = \E[-X] = -\E[X].
\end{equation}
\end{proof}

\section{More}

By Theorem~\ref{thm:main}, Definition~\ref{def:sphere} and \eqref{eq:sym}, see Section~\ref{sec:intro}.\footnote{A footnote.}
\begin{align}
  a &= b + c \label{eq:a}\\
  d &= e
\end{align}

\end{document}
